% Generated draft: requires mathematical and source review.
\documentclass[11pt]{article}
\usepackage[margin=25mm]{geometry}
\usepackage{fontspec}
\setmainfont{Latin Modern Roman}
\usepackage{amsmath,amssymb,mathtools,unicode-math}
\setmathfont{Latin Modern Math}
\usepackage{enumitem,needspace,xurl,hyperref}
\setlength{\parindent}{0pt}
\setlength{\parskip}{5pt}
\setlength{\emergencystretch}{3em}
\newcommand{\N}{\mathbb{N}}
\newcommand{\Z}{\mathbb{Z}}
\newcommand{\Q}{\mathbb{Q}}
\newcommand{\R}{\mathbb{R}}
\newcommand{\C}{\mathbb{C}}
\newcommand{\F}{\mathbb{F}}
\newcommand{\A}{\mathbb{A}}
\newcommand{\PP}{\mathbb{P}}
\newcommand{\E}{\mathbb{E}}
\newcommand{\eps}{\varepsilon}
\newcommand{\dd}{\,\mathrm{d}}
\begin{document}
\section*{UnsolvedMath ID: 701151}
\textbf{Source problem code:} catalogue entry
\subsection{Definitions and mathematical statement}
Let \(F_n=F(x_1,\ldots,x_n)\) be the free group of finite rank \(n\) with the displayed free basis. Define the semigroup
\[
S_n=\{x_{i_1}\cdots x_{i_k}:k\geq 1,\ 1\leq i_j\leq n\}.
\]
An element \(u\in F_n\) is \emph{positive} if \(u\in S_n\), and is \emph{potentially positive} if
\[
P_n(u)\quad\Longleftrightarrow\quad
\exists\alpha\in\operatorname{Aut}(F_n)\;\alpha(u)\in S_n.
\]
Here the semigroup convention uses nonempty positive words, so the identity is excluded. The source does not explicitly discuss the empty-word convention or rank zero.

\textbf{Part (a): algorithmic recognition.} Is potential positivity decidable? For a fixed finite rank \(n\), this asks whether there exists an algorithm which, on every input word in \(x_1^{\pm1},\ldots,x_n^{\pm1}\), halts and correctly decides whether its represented element satisfies \(P_n\). The source does not specify whether it requires a separate algorithm for each rank or one uniform algorithm taking both the rank and the word as input. This distinction remains unresolved in this draft.

The supplied comment dated 2013 reports that R. Goldstein answered part (a) affirmatively for \(n=2\), citing the 2006 reference listed below. This is a supplied literature report, not an independently verified assertion here.

\textbf{Part (b): archived question.} For \(m\geq n\), identify \(F_n\) with the subgroup generated by \(x_1,\ldots,x_n\) in \(F_m=F(x_1,\ldots,x_m)\), and write \(\iota_{n,m}:F_n\hookrightarrow F_m\) for this inclusion. This standard free-factor inclusion makes explicit the source's implicit interpretation of an element of \(F_n\) as an element of \(F_m\). Define stable potential positivity by
\[
\operatorname{SP}_n(u)\quad\Longleftrightarrow\quad
\exists m\geq n\;\exists\beta\in\operatorname{Aut}(F_m)\;
\beta(\iota_{n,m}(u))\in S_m.
\]
The archived question asks whether there exist a finite rank \(n\) and an element \(u\in F_n\) such that
\[
\operatorname{SP}_n(u)\ \land\ \neg P_n(u).
\]
The supplied archive reports a negative answer, attributed to A. Clark and R. Goldstein (2005); equivalently, it reports that stable potential positivity implies potential positivity in the original rank. This archival resolution is recorded as supplied attribution and is not independently verified here. The local record labels the overall problem \emph{partially solved}, retains part (a) in the main section, and places part (b) in the solved archive.
\subsection{Short English statement}
An element of a finitely generated free group is potentially positive if some automorphism turns it into a nonempty word using only the chosen generators, without inverses. Part (a) asks for an algorithm deciding this property; the supplied 2013 comment reports a positive answer in rank two. Part (b) asks whether adding free generators can make an element potentially positive when it was not so originally; the supplied archive reports that this cannot happen.
\subsection{Sources}
\begin{itemize}
\item Local UnsolvedMath registry, permanent ID \(701151\), exact source code \texttt{catalogue entry}, title \emph{Kourovka Notebook Problem 15.63}. Supplied upstream revision: \texttt{d371ee917ac9647fbb06cd6fc529deec6f3dda82}.
\item \emph{The Kourovka Notebook (Unsolved Problems in Group Theory)}, problem 15.63, proposed by A. G. Myasnikov and V. E. Shpilrain; 15th issue (2002), with the supplied record reporting an update in the 21st issue (2026), arXiv:1401.0300v48. The collection description names E. I. Khukhro and V. D. Mazurov as editors.
\item R. Goldstein, \emph{Contemp. Math.}, American Mathematical Society, volume 421 (2006), pages 157--168. The supplied comment dated 2013 attributes the affirmative rank-two answer to this reference. No article title is supplied.
\item A. Clark and R. Goldstein, \emph{Commun. Algebra}, volume 33, number 11 (2005), pages 4097--4104. The supplied Archive of Solved Problems entry attributes the negative answer to part (b) to this reference. No article title is supplied.
\item Supplied status review dated \texttt{2026-10-16}: canonical status \texttt{partially\_solved}, evidence marked \texttt{published}; its notes place solved part(s) in the archive and remaining part(s) in the main section. This is supplied dated curation, not an independently verified literature assessment. The registry creation and update date is also \texttt{2026-10-16}.
\end{itemize}
\end{document}
